% Generated by scripts/assemble.py; corrections belong in source/corrections.json.
\examhead{111}{安徽大学 2025 年数学分析试卷}{108}
% p109-b20; source PDF page 109
\par\Needspace{5\baselineskip}
\blockanchor{p109-b20}
\majorquestion{一}(16 分) 计算题.\par
% p109-b21; source PDF page 109
\blockanchor{p109-b21}
(1) (8 分) 设 $a_{n} = \sum_{k=1}^{n}\left(\sqrt{1 + \frac{k}{n^{2}}} - 1\right)$, 求 $\lim_{n \to \infty} a_{n}$.\par
% p109-b22; source PDF page 109
\blockanchor{p109-b22}
(2) (8 分) 求  $\lim_{x\to0^{+}}\frac{\int_{0}^{x}e^{-t^{2}}dt-x}{\sin x-x}.$\par
% p110-b00; source PDF page 110
\par\Needspace{5\baselineskip}
\blockanchor{p110-b00}
\majorquestion{二}（14分）证明 $f(x,y) = \left\{ \begin{array}{ll}\frac{xy}{\sqrt{x^2 + 4y^2}}, & x^2 +y^2\neq 0;\\ 0, & x^2 +y^2 = 0. \end{array} \right.$ 在原点 $(0,0)$ 处连续且可偏导，但不可微\par
% p110-b01; source PDF page 110
\par\Needspace{5\baselineskip}
\par\Needspace{6\baselineskip}
\blockanchor{p110-b01}
\majorquestion{三}(15 分) 计算\par
\nopagebreak[4]
% p110-b02; source PDF page 110
\blockanchor{p110-b02}
\displayformula{\iint_ {\Sigma} (x - y) \mathrm{d} y \mathrm{d} z + (y - z) \mathrm{d} z \mathrm{d} x + (z - x) \mathrm{d} x \mathrm{d} y.}
% p110-b03; source PDF page 110
\blockanchor{p110-b03}
其中  $\Sigma$  为曲面  $z = x^{2} + y^{2} (z \leq 1)$ ，方向向下.\par
% p110-b04; source PDF page 110
\par\Needspace{5\baselineskip}
\blockanchor{p110-b04}
\majorquestion{四}（15分）讨论反常积分 $\int_0^{+\infty}\frac{e^{\sin x}\sin 2x}{x^p}\mathrm{d}x$ 的敛散性，其中 $p > 0$\par
% p110-b05; source PDF page 110
\par\Needspace{5\baselineskip}
\blockanchor{p110-b05}
\majorquestion{五}(15 分) 设 $f(x) = \sum_{n=0}^{\infty} \frac{1}{2^n + x}$.\par
% p110-b06; source PDF page 110
\blockanchor{p110-b06}
(1) (10 分) 证明函数 $f(x)$ 在 $[0, +\infty)$ 上可导且一致连续.\par
% p110-b07; source PDF page 110
\blockanchor{p110-b07}
(2) (5 分) 证明反常积分 $\int_0^{+\infty} f(x) \, \mathrm{d}x$ 发散.\par
% p110-b08; source PDF page 110
\par\Needspace{5\baselineskip}
\par\Needspace{6\baselineskip}
\blockanchor{p110-b08}
\majorquestion{六}(15 分) 若函数 $f$ 一阶导数或偏导数存在且连续, 则记 $f \in C^{1}$. 设函数 $z = f(x, y)$ 由方程\par
\nopagebreak[4]
% p110-b09; source PDF page 110
\blockanchor{p110-b09}
\displayformula{F (x z - y, x - y z) = 0}
% p110-b10; source PDF page 110
\blockanchor{p110-b10}
唯一确定, 其中 $f \in C^{1}$, $F \in C^{1}$, 试求 $\oint_{L} (xz^{2} + 2yz) \mathrm{d}y - (2xz + yz^{2}) \mathrm{d}x$, 其中 $L$ 为正向单位圆周.\par
% p110-b11; source PDF page 110
\par\Needspace{5\baselineskip}
\blockanchor{p110-b11}
\majorquestion{七}(15 分) 若将 $a_1, a_2, \cdots, a_n, \cdots$ 重新排列后得到 $a_1', a_2', \cdots, a_n', \cdots$, 则级数 $\sum_{n=1}^{\infty} a_n'$ 称为 $\sum_{n=1}^{\infty} a_n$ 的一个更序级数. 设级数 $\sum_{n=1}^{\infty} x_n$ 条件收敛, 证明: 对任意给定的 $a, -\infty < a < +\infty$, 必定存在 $\sum_{n=1}^{\infty} x_n$ 的更序级数 $\sum_{n=1}^{\infty} x_n'$, 满足 $\sum_{n=1}^{\infty} x_n' = a$.\par
% p110-b12; source PDF page 110
\par\Needspace{5\baselineskip}
\blockanchor{p110-b12}
\majorquestion{八}(15 分) 设 $f(x)$ 在 $[0, +\infty)$ 上连续且有界, 若 $\forall r \in (-\infty, +\infty)$, $f(x) = r$ 在 $[0, +\infty)$ 上只有有限个根或无根, 证明: $\lim_{x \to +\infty} f(x)$ 存在.\par
% p110-b13; source PDF page 110
\par\Needspace{5\baselineskip}
\blockanchor{p110-b13}
\majorquestion{九}(15 分) 设有界函数 $f(x)$ 在闭区间 $[a, b]$ 上的不连续点为 $\{x_n\}_{n=1}^{\infty}$, 且 $\lim_{n \to \infty} x_n$ 存在, 证明: $f(x)$ 在闭区间 $[a, b]$ 上可积.\par
% p110-b14; source PDF page 110
\par\Needspace{5\baselineskip}
\par\Needspace{6\baselineskip}
\blockanchor{p110-b14}
\majorquestion{十}(15 分) 设 $f(x)$ 在闭区间 $[a, b]$ 上二阶可导, 且 $f(a) = f(b) = 0$. 证明: 对任意的 $x \in (a, b)$, 存在 $\xi \in (a, b)$, 使得\par
\nopagebreak[4]
% p110-b15; source PDF page 110
\blockanchor{p110-b15}
\displayformula{f (x) = \frac {f ^ {\prime \prime} (\xi) (x - a) (x - b)}{2}.}
